\documentclass[12pt,a4paper]{article}
\usepackage[utf8]{vietnam}
\usepackage{amsmath,amssymb}
\usepackage{tabularx}
\usepackage{tikz}
\usepackage{xcolor}
\usepackage[top=0.6cm, bottom=0.6cm, left=0.6cm, right=0.6cm]{geometry}
\usepackage[dethi]{ex_test}

\SAOPTN{widthSA=2.8,heightSA=0.62}
\setlength{\parskip}{0pt}

% Dòng kẻ chấm mờ chuẩn viết tay
\newcommand{\dongke}[1]{%
  \begingroup%
  \setlength{\parskip}{0pt}%
  \baselineskip=18.5pt%
  \vspace{1.0mm}%
  \foreach \i in {1,...,#1}{%
    \noindent{\color{gray!60}\dotfill}\par%
  }%
  \vspace{0.6mm}%
  \endgroup%
}

% Định nghĩa ô bài tập tự luận trong bảng
\newcommand{\oBaiTap}[3][6]{%
  \vspace{1.0mm}%
  \noindent\textbf{#2} #3\par%
  \dongke{#1}%
}

\begin{document}
\pagestyle{empty}
% =============================================================================
% ========================= MA DE 101 - MAT TRUOC (TRANG 1/2) =============
% =============================================================================
\noindent
\begin{tabularx}{\linewidth}{|X|c|}
\hline
\textbf{PHIẾU KIỂM TRA ĐỊNH KỲ TOÁN 6 - TUẦN 11 - BUỔI 1} & \textbf{Mã đề: 101 - Trang 1/2} \\
\hline
\end{tabularx}

\vspace{1.0mm}

\noindent
\begin{tabularx}{\linewidth}{|p{3.2cm}|X|}
\hline
\centering\textbf{Điểm số} & \centering\textbf{Lời nhận xét của thầy cô} \tabularnewline
\rule{0pt}{1.8cm} & \\
\hline
\end{tabularx}

\vspace{1.0mm}
\noindent\textbf{Câu 1 (0,25 điểm).} Kết quả của phép tính $(-15) + (-25)$ là: \shortans{-40}\par\vspace{0.25mm}
\noindent\textbf{Câu 2 (0,25 điểm).} Kết quả của phép tính $45 + (-20)$ là: \shortans{25}\par\vspace{0.25mm}
\noindent\textbf{Câu 3 (0,25 điểm).} Kết quả của phép tính $(-65) + 30$ là: \shortans{-35}\par\vspace{0.25mm}
\noindent\textbf{Câu 4 (0,25 điểm).} Kết quả của phép tính $(-2026) + 2026$ là: \shortans{0}\par\vspace{0.25mm}
\noindent\textbf{Câu 5 (0,25 điểm).} Số nguyên thích hợp điền vào chỗ chấm $(-37) + \dots = 0$ là: \shortans{37}\par\vspace{0.25mm}
\noindent\textbf{Câu 6 (0,25 điểm).} Điền dấu thích hợp ($<, >, =$) vào chỗ chấm: $(-12) + (-18) \; \dots \; 0$ \shortans{<}\par\vspace{0.25mm}
\noindent\textbf{Câu 7 (0,25 điểm).} Trong các chữ cái N, F, A, G, chữ cái có trục đối xứng là: \shortans{A}\par\vspace{0.25mm}
\noindent\textbf{Câu 8 (0,25 điểm).} Trong các chữ cái A, S, M, E, chữ cái có tâm đối xứng là: \shortans{S}\par\vspace{0.25mm}
\noindent\textbf{Câu 9 (0,25 điểm).} Một hình thoi có độ dài cạnh $8\text{ cm}$. Chu vi hình thoi đó là: \shortans{32\text{ cm}}\par\vspace{0.25mm}
\noindent\textbf{Câu 10 (0,25 điểm).} Nhiệt độ Sa Pa lúc sáng $-2^\circ\text{C}$, trưa tăng $7^\circ\text{C}$. Nhiệt độ trưa là: \shortans{5^\circ\text{C}}\par\vspace{0.25mm}

\vspace{0.3mm}
\noindent\textbf{Câu 11 (3,0 điểm). Thực hiện phép tính (hợp lý nếu có thể):}
\vspace{0.4mm}

\noindent
\begin{tabularx}{\linewidth}{|X|X|}
\hline
\oBaiTap[7]{a)}{$(-124) + (-76)$} & \oBaiTap[7]{b)}{$(-165) + 215$} \\
\hline
\oBaiTap[7]{c)}{$145 + (-68) + (-145) + 68$} & \oBaiTap[7]{d)}{$(-45) + 183 + (-55) + 17$} \\
\hline
\oBaiTap[7]{e)}{$(-125) + 25 \cdot 2^2$} & \oBaiTap[7]{f)}{$(-342) + [342 + (-158)]$} \\
\hline
\end{tabularx}
\newpage
% =============================================================================
% ========================= MA DE 101 - MAT SAU (TRANG 2/2) ==============
% =============================================================================
\noindent
\begin{tabularx}{\linewidth}{|X|}
\hline
\centering\textbf{TOÁN 6 - TUẦN 11 - BUỔI 1 | Mã đề: 101 - Trang 2/2} \\
\hline
\end{tabularx}

\vspace{1.2mm}

\noindent\textbf{Câu 12 (2,0 điểm). Tìm số nguyên $x$, biết:}
\vspace{0.5mm}

\noindent
\begin{tabularx}{\linewidth}{|X|X|}
\hline
\oBaiTap[5]{a)}{$x + 18 = 10$} & \oBaiTap[5]{b)}{$x + (-25) = -15$} \\
\hline
\oBaiTap[6]{c)}{$3x + (-14) = 16$} & \oBaiTap[6]{d)}{Liệt kê và tính tổng các số nguyên $x$ thỏa mãn: $-5 \le x < 6$} \\
\hline
\end{tabularx}

\vspace{1.2mm}

\noindent
\begin{tabularx}{\linewidth}{|X|}
\hline
\vspace{0.4mm}
\textbf{Câu 13 (2,0 điểm).} Đầu ngày, số dư tài khoản ngân hàng của một cửa hàng là $50\text{ triệu đồng}$. Trong ngày, cửa hàng thực hiện 4 giao dịch (tiền thu vào ghi dấu $+$, tiền chi ra ghi dấu $-$) được ghi nhận trong bảng sau:\par\vspace{0.4mm}\centerline{\begin{tabular}{|l|c|c|c|c|}\hline \textbf{Giao dịch} & GD 1 (Thu tiền hàng) & GD 2 (Nhập hàng mới) & GD 3 (Thu dịch vụ) & GD 4 (Trả tiền thuê) \\\hline \textbf{Số tiền} & $+35\text{ triệu đồng}$ & $-40\text{ triệu đồng}$ & $+25\text{ triệu đồng}$ & $-30\text{ triệu đồng}$ \\\hline\end{tabular}}\vspace{0.4mm}\par\vspace{0.3mm}
\textbf{a) (1,0 điểm)} Tính tổng số tiền thu vào và tổng số tiền chi ra của cửa hàng trong ngày hôm đó.\par\vspace{0.3mm}
\textbf{b) (1,0 điểm)} Tính số dư tài khoản của cửa hàng vào cuối ngày hôm đó.\par
\vspace{0.4mm}
\noindent\textit{Bài giải:}\dongke{10}
\\ \hline
\end{tabularx}

\vspace{1.2mm}

\noindent
\begin{tabularx}{\linewidth}{|X|}
\hline
\vspace{0.4mm}
\textbf{Câu 14 (0,5 điểm).} Tính giá trị của biểu thức: $S = 1 + (-2) + 3 + (-4) + 5 + (-6) + \dots + 2025 + (-2026)$.\par
\vspace{0.4mm}
\noindent\textit{Bài giải:}\dongke{7}
\\ \hline
\end{tabularx}

\newpage
% =============================================================================
% ========================= MA DE 102 - MAT TRUOC (TRANG 1/2) =============
% =============================================================================
\noindent
\begin{tabularx}{\linewidth}{|X|c|}
\hline
\textbf{PHIẾU KIỂM TRA ĐỊNH KỲ TOÁN 6 - TUẦN 11 - BUỔI 1} & \textbf{Mã đề: 102 - Trang 1/2} \\
\hline
\end{tabularx}

\vspace{1.0mm}

\noindent
\begin{tabularx}{\linewidth}{|p{3.2cm}|X|}
\hline
\centering\textbf{Điểm số} & \centering\textbf{Lời nhận xét của thầy cô} \tabularnewline
\rule{0pt}{1.8cm} & \\
\hline
\end{tabularx}

\vspace{1.0mm}
\noindent\textbf{Câu 1 (0,25 điểm).} Kết quả của phép tính $(-18) + (-22)$ là: \shortans{-40}\par\vspace{0.25mm}
\noindent\textbf{Câu 2 (0,25 điểm).} Kết quả của phép tính $58 + (-30)$ là: \shortans{28}\par\vspace{0.25mm}
\noindent\textbf{Câu 3 (0,25 điểm).} Kết quả của phép tính $(-72) + 22$ là: \shortans{-50}\par\vspace{0.25mm}
\noindent\textbf{Câu 4 (0,25 điểm).} Kết quả của phép tính $2025 + (-2025)$ là: \shortans{0}\par\vspace{0.25mm}
\noindent\textbf{Câu 5 (0,25 điểm).} Số nguyên thích hợp điền vào chỗ chấm $\dots + 48 = 0$ là: \shortans{-48}\par\vspace{0.25mm}
\noindent\textbf{Câu 6 (0,25 điểm).} Điền dấu thích hợp ($<, >, =$) vào chỗ chấm: $35 + (-15) \; \dots \; 0$ \shortans{>}\par\vspace{0.25mm}
\noindent\textbf{Câu 7 (0,25 điểm).} Trong các chữ cái P, M, R, S, chữ cái có trục đối xứng là: \shortans{M}\par\vspace{0.25mm}
\noindent\textbf{Câu 8 (0,25 điểm).} Trong các chữ cái B, Z, P, C, chữ cái có tâm đối xứng là: \shortans{Z}\par\vspace{0.25mm}
\noindent\textbf{Câu 9 (0,25 điểm).} Một hình thoi có độ dài cạnh $12\text{ cm}$. Chu vi hình thoi đó là: \shortans{48\text{ cm}}\par\vspace{0.25mm}
\noindent\textbf{Câu 10 (0,25 điểm).} Nhiệt độ Mẫu Sơn lúc sáng $-4^\circ\text{C}$, trưa tăng $8^\circ\text{C}$. Nhiệt độ trưa là: \shortans{4^\circ\text{C}}\par\vspace{0.25mm}

\vspace{0.3mm}
\noindent\textbf{Câu 11 (3,0 điểm). Thực hiện phép tính (hợp lý nếu có thể):}
\vspace{0.4mm}

\noindent
\begin{tabularx}{\linewidth}{|X|X|}
\hline
\oBaiTap[7]{a)}{$(-135) + (-65)$} & \oBaiTap[7]{b)}{$(-175) + 235$} \\
\hline
\oBaiTap[7]{c)}{$156 + (-74) + (-156) + 74$} & \oBaiTap[7]{d)}{$(-38) + 164 + (-62) + 36$} \\
\hline
\oBaiTap[7]{e)}{$(-150) + 30 \cdot 2^2$} & \oBaiTap[7]{f)}{$(-425) + [425 + (-172)]$} \\
\hline
\end{tabularx}
\newpage
% =============================================================================
% ========================= MA DE 102 - MAT SAU (TRANG 2/2) ==============
% =============================================================================
\noindent
\begin{tabularx}{\linewidth}{|X|}
\hline
\centering\textbf{TOÁN 6 - TUẦN 11 - BUỔI 1 | Mã đề: 102 - Trang 2/2} \\
\hline
\end{tabularx}

\vspace{1.2mm}

\noindent\textbf{Câu 12 (2,0 điểm). Tìm số nguyên $x$, biết:}
\vspace{0.5mm}

\noindent
\begin{tabularx}{\linewidth}{|X|X|}
\hline
\oBaiTap[5]{a)}{$x + 24 = 14$} & \oBaiTap[5]{b)}{$x + (-30) = -18$} \\
\hline
\oBaiTap[6]{c)}{$2x + (-16) = 24$} & \oBaiTap[6]{d)}{Liệt kê và tính tổng các số nguyên $x$ thỏa mãn: $-6 < x \le 5$} \\
\hline
\end{tabularx}

\vspace{1.2mm}

\noindent
\begin{tabularx}{\linewidth}{|X|}
\hline
\vspace{0.4mm}
\textbf{Câu 13 (2,0 điểm).} Đầu ngày, số dư tài khoản ngân hàng của một cửa hàng là $60\text{ triệu đồng}$. Trong ngày, cửa hàng thực hiện 4 giao dịch (tiền thu vào ghi dấu $+$, tiền chi ra ghi dấu $-$) được ghi nhận trong bảng sau:\par\vspace{0.4mm}\centerline{\begin{tabular}{|l|c|c|c|c|}\hline \textbf{Giao dịch} & GD 1 (Thu tiền hàng) & GD 2 (Nhập hàng mới) & GD 3 (Thu dịch vụ) & GD 4 (Trả tiền thuê) \\\hline \textbf{Số tiền} & $+45\text{ triệu đồng}$ & $-50\text{ triệu đồng}$ & $+30\text{ triệu đồng}$ & $-35\text{ triệu đồng}$ \\\hline\end{tabular}}\vspace{0.4mm}\par\vspace{0.3mm}
\textbf{a) (1,0 điểm)} Tính tổng số tiền thu vào và tổng số tiền chi ra của cửa hàng trong ngày hôm đó.\par\vspace{0.3mm}
\textbf{b) (1,0 điểm)} Tính số dư tài khoản của cửa hàng vào cuối ngày hôm đó.\par
\vspace{0.4mm}
\noindent\textit{Bài giải:}\dongke{10}
\\ \hline
\end{tabularx}

\vspace{1.2mm}

\noindent
\begin{tabularx}{\linewidth}{|X|}
\hline
\vspace{0.4mm}
\textbf{Câu 14 (0,5 điểm).} Tính giá trị của biểu thức: $S = 1 + (-2) + 3 + (-4) + 5 + (-6) + \dots + 2023 + (-2024)$.\par
\vspace{0.4mm}
\noindent\textit{Bài giải:}\dongke{7}
\\ \hline
\end{tabularx}

\newpage
% =============================================================================
% ========================= MA DE 103 - MAT TRUOC (TRANG 1/2) =============
% =============================================================================
\noindent
\begin{tabularx}{\linewidth}{|X|c|}
\hline
\textbf{PHIẾU KIỂM TRA ĐỊNH KỲ TOÁN 6 - TUẦN 11 - BUỔI 1} & \textbf{Mã đề: 103 - Trang 1/2} \\
\hline
\end{tabularx}

\vspace{1.0mm}

\noindent
\begin{tabularx}{\linewidth}{|p{3.2cm}|X|}
\hline
\centering\textbf{Điểm số} & \centering\textbf{Lời nhận xét của thầy cô} \tabularnewline
\rule{0pt}{1.8cm} & \\
\hline
\end{tabularx}

\vspace{1.0mm}
\noindent\textbf{Câu 1 (0,25 điểm).} Kết quả của phép tính $(-16) + (-34)$ là: \shortans{-50}\par\vspace{0.25mm}
\noindent\textbf{Câu 2 (0,25 điểm).} Kết quả của phép tính $62 + (-22)$ là: \shortans{40}\par\vspace{0.25mm}
\noindent\textbf{Câu 3 (0,25 điểm).} Kết quả của phép tính $(-84) + 34$ là: \shortans{-50}\par\vspace{0.25mm}
\noindent\textbf{Câu 4 (0,25 điểm).} Kết quả của phép tính $(-1985) + 1985$ là: \shortans{0}\par\vspace{0.25mm}
\noindent\textbf{Câu 5 (0,25 điểm).} Số nguyên thích hợp điền vào chỗ chấm $(-56) + \dots = 0$ là: \shortans{56}\par\vspace{0.25mm}
\noindent\textbf{Câu 6 (0,25 điểm).} Điền dấu thích hợp ($<, >, =$) vào chỗ chấm: $(-45) + (-25) \; \dots \; 0$ \shortans{<}\par\vspace{0.25mm}
\noindent\textbf{Câu 7 (0,25 điểm).} Trong các chữ cái Q, Y, L, J, chữ cái có trục đối xứng là: \shortans{Y}\par\vspace{0.25mm}
\noindent\textbf{Câu 8 (0,25 điểm).} Trong các chữ cái N, D, F, G, chữ cái có tâm đối xứng là: \shortans{N}\par\vspace{0.25mm}
\noindent\textbf{Câu 9 (0,25 điểm).} Một hình thoi có độ dài cạnh $15\text{ cm}$. Chu vi hình thoi đó là: \shortans{60\text{ cm}}\par\vspace{0.25mm}
\noindent\textbf{Câu 10 (0,25 điểm).} Tàu ngầm ở độ sâu $-35\text{ m}$, sau đó nổi lên thêm $15\text{ m}$. Độ sâu mới là: \shortans{-20\text{ m}}\par\vspace{0.25mm}

\vspace{0.3mm}
\noindent\textbf{Câu 11 (3,0 điểm). Thực hiện phép tính (hợp lý nếu có thể):}
\vspace{0.4mm}

\noindent
\begin{tabularx}{\linewidth}{|X|X|}
\hline
\oBaiTap[7]{a)}{$(-148) + (-52)$} & \oBaiTap[7]{b)}{$(-185) + 255$} \\
\hline
\oBaiTap[7]{c)}{$168 + (-83) + (-168) + 83$} & \oBaiTap[7]{d)}{$(-47) + 172 + (-53) + 28$} \\
\hline
\oBaiTap[7]{e)}{$(-180) + 35 \cdot 2^2$} & \oBaiTap[7]{f)}{$(-516) + [516 + (-184)]$} \\
\hline
\end{tabularx}
\newpage
% =============================================================================
% ========================= MA DE 103 - MAT SAU (TRANG 2/2) ==============
% =============================================================================
\noindent
\begin{tabularx}{\linewidth}{|X|}
\hline
\centering\textbf{TOÁN 6 - TUẦN 11 - BUỔI 1 | Mã đề: 103 - Trang 2/2} \\
\hline
\end{tabularx}

\vspace{1.2mm}

\noindent\textbf{Câu 12 (2,0 điểm). Tìm số nguyên $x$, biết:}
\vspace{0.5mm}

\noindent
\begin{tabularx}{\linewidth}{|X|X|}
\hline
\oBaiTap[5]{a)}{$x + 32 = 20$} & \oBaiTap[5]{b)}{$x + (-35) = -20$} \\
\hline
\oBaiTap[6]{c)}{$4x + (-18) = 22$} & \oBaiTap[6]{d)}{Liệt kê và tính tổng các số nguyên $x$ thỏa mãn: $-4 \le x < 5$} \\
\hline
\end{tabularx}

\vspace{1.2mm}

\noindent
\begin{tabularx}{\linewidth}{|X|}
\hline
\vspace{0.4mm}
\textbf{Câu 13 (2,0 điểm).} Đầu ngày, số dư tài khoản ngân hàng của một cửa hàng là $70\text{ triệu đồng}$. Trong ngày, cửa hàng thực hiện 4 giao dịch (tiền thu vào ghi dấu $+$, tiền chi ra ghi dấu $-$) được ghi nhận trong bảng sau:\par\vspace{0.4mm}\centerline{\begin{tabular}{|l|c|c|c|c|}\hline \textbf{Giao dịch} & GD 1 (Thu tiền hàng) & GD 2 (Nhập hàng mới) & GD 3 (Thu dịch vụ) & GD 4 (Trả tiền thuê) \\\hline \textbf{Số tiền} & $+40\text{ triệu đồng}$ & $-60\text{ triệu đồng}$ & $+35\text{ triệu đồng}$ & $-45\text{ triệu đồng}$ \\\hline\end{tabular}}\vspace{0.4mm}\par\vspace{0.3mm}
\textbf{a) (1,0 điểm)} Tính tổng số tiền thu vào và tổng số tiền chi ra của cửa hàng trong ngày hôm đó.\par\vspace{0.3mm}
\textbf{b) (1,0 điểm)} Tính số dư tài khoản của cửa hàng vào cuối ngày hôm đó.\par
\vspace{0.4mm}
\noindent\textit{Bài giải:}\dongke{10}
\\ \hline
\end{tabularx}

\vspace{1.2mm}

\noindent
\begin{tabularx}{\linewidth}{|X|}
\hline
\vspace{0.4mm}
\textbf{Câu 14 (0,5 điểm).} Tính giá trị của biểu thức: $S = 1 + (-2) + 3 + (-4) + 5 + (-6) + \dots + 2027 + (-2028)$.\par
\vspace{0.4mm}
\noindent\textit{Bài giải:}\dongke{7}
\\ \hline
\end{tabularx}

\newpage
% =============================================================================
% ========================= MA DE 104 - MAT TRUOC (TRANG 1/2) =============
% =============================================================================
\noindent
\begin{tabularx}{\linewidth}{|X|c|}
\hline
\textbf{PHIẾU KIỂM TRA ĐỊNH KỲ TOÁN 6 - TUẦN 11 - BUỔI 1} & \textbf{Mã đề: 104 - Trang 1/2} \\
\hline
\end{tabularx}

\vspace{1.0mm}

\noindent
\begin{tabularx}{\linewidth}{|p{3.2cm}|X|}
\hline
\centering\textbf{Điểm số} & \centering\textbf{Lời nhận xét của thầy cô} \tabularnewline
\rule{0pt}{1.8cm} & \\
\hline
\end{tabularx}

\vspace{1.0mm}
\noindent\textbf{Câu 1 (0,25 điểm).} Kết quả của phép tính $(-27) + (-23)$ là: \shortans{-50}\par\vspace{0.25mm}
\noindent\textbf{Câu 2 (0,25 điểm).} Kết quả của phép tính $75 + (-35)$ là: \shortans{40}\par\vspace{0.25mm}
\noindent\textbf{Câu 3 (0,25 điểm).} Kết quả của phép tính $(-95) + 45$ là: \shortans{-50}\par\vspace{0.25mm}
\noindent\textbf{Câu 4 (0,25 điểm).} Kết quả của phép tính $1994 + (-1994)$ là: \shortans{0}\par\vspace{0.25mm}
\noindent\textbf{Câu 5 (0,25 điểm).} Số nguyên thích hợp điền vào chỗ chấm $\dots + 65 = 0$ là: \shortans{-65}\par\vspace{0.25mm}
\noindent\textbf{Câu 6 (0,25 điểm).} Điền dấu thích hợp ($<, >, =$) vào chỗ chấm: $60 + (-20) \; \dots \; 0$ \shortans{>}\par\vspace{0.25mm}
\noindent\textbf{Câu 7 (0,25 điểm).} Trong các chữ cái K, T, G, R, chữ cái có trục đối xứng là: \shortans{T}\par\vspace{0.25mm}
\noindent\textbf{Câu 8 (0,25 điểm).} Trong các chữ cái H, V, W, Y, chữ cái có tâm đối xứng là: \shortans{H}\par\vspace{0.25mm}
\noindent\textbf{Câu 9 (0,25 điểm).} Một hình thoi có độ dài cạnh $25\text{ cm}$. Chu vi hình thoi đó là: \shortans{100\text{ cm}}\par\vspace{0.25mm}
\noindent\textbf{Câu 10 (0,25 điểm).} Tàu ngầm ở độ sâu $-40\text{ m}$, sau đó nổi lên thêm $25\text{ m}$. Độ sâu mới là: \shortans{-15\text{ m}}\par\vspace{0.25mm}

\vspace{0.3mm}
\noindent\textbf{Câu 11 (3,0 điểm). Thực hiện phép tính (hợp lý nếu có thể):}
\vspace{0.4mm}

\noindent
\begin{tabularx}{\linewidth}{|X|X|}
\hline
\oBaiTap[7]{a)}{$(-156) + (-44)$} & \oBaiTap[7]{b)}{$(-195) + 275$} \\
\hline
\oBaiTap[7]{c)}{$172 + (-95) + (-172) + 95$} & \oBaiTap[7]{d)}{$(-29) + 158 + (-71) + 42$} \\
\hline
\oBaiTap[7]{e)}{$(-200) + 40 \cdot 2^2$} & \oBaiTap[7]{f)}{$(-635) + [635 + (-196)]$} \\
\hline
\end{tabularx}
\newpage
% =============================================================================
% ========================= MA DE 104 - MAT SAU (TRANG 2/2) ==============
% =============================================================================
\noindent
\begin{tabularx}{\linewidth}{|X|}
\hline
\centering\textbf{TOÁN 6 - TUẦN 11 - BUỔI 1 | Mã đề: 104 - Trang 2/2} \\
\hline
\end{tabularx}

\vspace{1.2mm}

\noindent\textbf{Câu 12 (2,0 điểm). Tìm số nguyên $x$, biết:}
\vspace{0.5mm}

\noindent
\begin{tabularx}{\linewidth}{|X|X|}
\hline
\oBaiTap[5]{a)}{$x + 28 = 16$} & \oBaiTap[5]{b)}{$x + (-40) = -25$} \\
\hline
\oBaiTap[6]{c)}{$3x + (-15) = 18$} & \oBaiTap[6]{d)}{Liệt kê và tính tổng các số nguyên $x$ thỏa mãn: $-5 < x \le 4$} \\
\hline
\end{tabularx}

\vspace{1.2mm}

\noindent
\begin{tabularx}{\linewidth}{|X|}
\hline
\vspace{0.4mm}
\textbf{Câu 13 (2,0 điểm).} Đầu ngày, số dư tài khoản ngân hàng của một cửa hàng là $80\text{ triệu đồng}$. Trong ngày, cửa hàng thực hiện 4 giao dịch (tiền thu vào ghi dấu $+$, tiền chi ra ghi dấu $-$) được ghi nhận trong bảng sau:\par\vspace{0.4mm}\centerline{\begin{tabular}{|l|c|c|c|c|}\hline \textbf{Giao dịch} & GD 1 (Thu tiền hàng) & GD 2 (Nhập hàng mới) & GD 3 (Thu dịch vụ) & GD 4 (Trả tiền thuê) \\\hline \textbf{Số tiền} & $+50\text{ triệu đồng}$ & $-65\text{ triệu đồng}$ & $+40\text{ triệu đồng}$ & $-45\text{ triệu đồng}$ \\\hline\end{tabular}}\vspace{0.4mm}\par\vspace{0.3mm}
\textbf{a) (1,0 điểm)} Tính tổng số tiền thu vào và tổng số tiền chi ra của cửa hàng trong ngày hôm đó.\par\vspace{0.3mm}
\textbf{b) (1,0 điểm)} Tính số dư tài khoản của cửa hàng vào cuối ngày hôm đó.\par
\vspace{0.4mm}
\noindent\textit{Bài giải:}\dongke{10}
\\ \hline
\end{tabularx}

\vspace{1.2mm}

\noindent
\begin{tabularx}{\linewidth}{|X|}
\hline
\vspace{0.4mm}
\textbf{Câu 14 (0,5 điểm).} Tính giá trị của biểu thức: $S = 1 + (-2) + 3 + (-4) + 5 + (-6) + \dots + 2021 + (-2022)$.\par
\vspace{0.4mm}
\noindent\textit{Bài giải:}\dongke{7}
\\ \hline
\end{tabularx}
\end{document}
